\subsection{有限域的代数扩张}

命题 3. --- 设 K 为含 q 个元素的有限域，Ω 为 K 的代数闭扩张，且 m 为整数 ≥1。

\begin{enumerate}
\def\labelenumi{\alph{enumi})}
\item
  存在唯一的子扩张 \(\mathbf{K}_m\) （ \(\Omega\) 的），其次数为 \(m\) （在 \(\mathbf{K}\) 上）。
\item
  域 \(K_{m}\) 具有 \(q^{m}\) 个元素，并且是自同构 \(x \mapsto x^{q^{m}}\) （ \(\Omega\) 的）的不动点集。
\item
  我们有 \(\mathbf{K}_m = \mathbf{K}(\zeta)\) ，其中 \(\zeta\) 是循环群 \(\mathbf{K}_m^*\) 的一个生成元。
\end{enumerate}

设 \(p\) 为 \(\mathbf{K}\) 的特征， \(f\) 为 \(\mathbf{K}\) 在 \(\mathbf{F}_p\) 上的次数。我们有 \(q^m = p^{fm}\) ，且映射 \(x \mapsto x^{q^m}\) 因而是完全域 \(\Omega\) 的自同构（V，第7页，命题4）。因此， \(\mathbf{K}_m\) （即多项式 \(X^{q^m} - X\) （ \(\mathbf{K}[\mathbf{X}]\) 中的）的根集）是 \(\Omega\) 的一个子域。由于 \(X^{q^m} - X\) 的导数等于 -1，该多项式的所有根都是单根（IV，第17页，命题7），所以 \(\mathbf{K}_m\) 具有 \(q^m\) 个元素。由此可得 \([\mathbf{K}_m : \mathbf{K}] = m\) 。

现在设 L 是 \(\Omega\) 的一个子扩张，在 K 上的次数为 \(m\) 。作为 K 上的向量空间，L 同构于 \(\mathbf{K}^m\) ，因此有 \(q^m\) 个元素。于是我们有 \(x^{q^m} = x\) 对所有 \(x \in L\) 成立（命题 2），从而 \(\mathrm{L} \subset \mathrm{K}_m\) 。由于 \([\mathrm{L} : \mathrm{K}] = [\mathrm{K}_m : \mathrm{K}] = m\) ，我们最终得到 \(\mathrm{L} = \mathrm{K}_m\) 。

这样我们就证明了断言 \(a)\) 与 \(b)\) ，而 \(c)\) 是平凡的。

推论. --- 设 K 为有限域，且 \(\Omega\) 为 K 的一个代数闭扩张。相对代数闭包 \(\bar{K}\) （ K 在 \(\Omega\) 中的）由 0 和单位根组成，并且是 K 的一个代数闭包。

我们知道 \(\bar{K}\) 是 K 的一个代数闭包（V，第22页，例2），并且显然 \(\Omega\) 中的每一个单位根都属于 \(\bar{K}\) 。进一步地，给定 \(x \neq 0\) （ \(\bar{K}\) 中的），其在 K 上次数为 m；若 K 有 q 个元素，则域 \(\mathbf{K}(x)\) 具有 \(q^{m}\) 个元素，从而 \(x^{q^{m}-1} = 1\) ，因此 x 是 \(\Omega\) 中的一个单位根。

设 \(p\) 为一个素数，并设 \(\mathbf{F}_p = \mathbf{Z} / p\mathbf{Z}\) 为 \(p\) 个元素的域。我们选取一个代数闭包 \(\Omega\) （ \(\mathbf{F}_p\) 的），其存在性由施泰尼茨定理（V，第23页，定理2）保证。设 \(f\) 为正整数且 \(q = p^f\) 。根据命题 3，存在 \(\Omega\) 的唯一子域，其次数为 \(f\) （在 \(\mathbf{F}_p\) 上）；我们将其记为 \(\mathbf{F}_q(\Omega)\) ，或（在滥用记号的意义下）记为 \(\mathbf{F}_q\) 。它是 \(\Omega\) 的唯一子扩张，其次数为 \(f\) （在 \(\mathbf{F}_p\) 上）。它是 \(\Omega\) 的唯一子域，基数为 \(q\) ，并且每个基数为 \(q\) 的域都（非典范地）同构于 \(\mathbf{F}_q\) （命题 2 的推论）。我们注意到 \(\mathbf{F}_q\) 由那些 \(x\) （ \(\Omega\) 中的、使得 \(x^q = x\) 的）组成，且 \(\mathbf{F}_q \subset \mathbf{F}_{q'}\) 成立，当且仅当 \(q'\) 是 \(q\) 的幂。

命题 4. --- 设 K 是含 q 个元素的有限域，且 \(K_{m}\) 是 K 的一个有限次扩张，次数为 m。

\begin{enumerate}
\def\labelenumi{\alph{enumi})}
\item
  域 \(K_{m}\) 是 K 的伽罗瓦扩张，其伽罗瓦群是由自同构 \(\sigma_{q}: x \mapsto x^{q}\) 生成的 m 阶循环群。
\item
  对每个 \(x \in \mathbf{K}_m\) ， \(x\) 关于 \(\mathbf{K}\) 的范数等于 \(x^{(q^m - 1) / (q - 1)}\) 。
\item
  \(\mathbf{K}\) 的每个元素都是 \(\mathbf{K}_m\) 的某个元素的迹（相应地，范数）。
\end{enumerate}

设 \(\Gamma\) 为 \(\mathbf{K}_m\) 的由 \(\sigma_q\) 生成的自同构循环群。 \(\Gamma\) 的不变域由元素 \(x\) （属于 \(\mathbf{K}_m\) ，使得 \(x^q = x\) ）组成，因此等于 \(\mathbf{K}\) 。所以 \(\mathbf{K}_m\) 是 \(\mathbf{K}\) 的伽罗瓦扩张，其伽罗瓦群为 \(\Gamma\) ，且后者的阶等于 \([\mathbf{K}_m: \mathbf{K}] = m\) （V，第66页，定理3）。因此 \(a)\) 由此得出。

我们有 \(\Gamma = \{1, \sigma_q, \sigma_q^2, \ldots, \sigma_q^{m-1}\}\) ；一个元素 \(x\) （ \(\mathbf{K}_m\) 中的）关于 \(\mathbf{K}\) 的范数因而是 \(\mathbf{N}(x) = \prod_{i=0}^{m-1} \sigma_q^i(x) = x^{1+q+\cdots+q^{m-1}}\) ，并且我们有 \(1 + q + \cdots + q^{m-1} = \frac{q^m - 1}{q - 1}\) 。这证明了 \(b\) ）。设 \(\zeta\) 是循环群 \(\mathbf{K}_m^*\) 的一个生成元；范数 \(\mathbf{N}: \mathbf{K}_M^* \to \mathbf{K}^*\) 的像是 \(\mathbf{K}_i^*\) 的由元素 \(\xi = \mathbf{N}(\zeta) = \zeta^{(q^m - 1)/(q - 1)}\) 生成的循环子群；由于 \(\zeta\) 的阶为 \(q^m - 1\) ， \(\xi\) 的阶为 \(q - 1\) ，因而生成 \(\mathbf{K}^*\) 。这证明了 K 的每个非零元素都是 \(K_{m}\) 的某个非零元素的范数；此外，我们有 \(0 = \mathbf{N}(0)\) 。

最后，由于 \(K_{m}\) 是 K 的可分代数扩张，迹是向量空间 \(K_{m}\) （K 上的）的非零线性形式（V，第49页，推论）；因此 K 的每个元素都是 \(K_{m}\) 的某个元素的迹。
% semantic-fragments: legacy-input-seam
\relax{}
